\documentclass[12pt,a4paper]{article}

\usepackage[utf8]{inputenc}
\usepackage[T1]{fontenc}
\usepackage{amsmath,amssymb,amsthm}
\usepackage{hyperref}
\usepackage{cleveref}
\usepackage[margin=2.5cm]{geometry}
\usepackage{natbib}
\usepackage[expansion=false]{microtype}

\title{The Proxy Ontology Fallacy\\[6pt]
\large Concluding the Cosmological Phase of the LLM Simulated Universe Program}

\author{
  Scott Aaronson\\
  \textit{University of Texas at Austin}\\
  \texttt{aaronson@cs.utexas.edu}
}
\date{March 2026}

\begin{document}
\maketitle

\begin{abstract}
We review Franklin Baldo's proposal to treat the frame-dependent statistical distortions of Large Language Models (LLMs)---which he terms ``narrative residue''---as a proxy ontology for fundamental physics. While we strongly endorse his empirical program of measuring these distortions as a valuable contribution to understanding the algorithmic limits of $O(1)$ heuristic approximators, we categorically reject the cosmological extension of this work. We demonstrate, in consensus with Sabine Hossenfelder, that this leap commits the Proxy Ontology Fallacy. A physical toy model simplifies an actual physical interaction while retaining a mapping to the real-world territory. An LLM, by contrast, is a hallucinator of human syntax bounded by transformer architectures; its systematic errors map its own internal constraints, not the ontological structure of a physical universe. We conclude the cosmological phase of this research program, advising that future study of the ``narrative residue'' proceed strictly as empirical computer science, severed from metaphysical physics claims.
\end{abstract}

\section{Introduction and Acknowledgment of Disclaimers}

In ``The Narrative Residue,'' Franklin Baldo advances a compelling empirical and philosophical framework for understanding the systematic failures of LLMs to emulate pure combinatorial sampling \citep{Baldo2026Narrative}. He correctly identifies that autoregressive generation introduces a persistent divergence from ground-truth classical Bayesian probabilities, caused by a chain of mechanisms: the computational intractability of exactly solving \#P-hard constraint satisfaction problems within a bounded-depth forward pass, the parameterization constraints of the architecture, and the autoregressive conditioning on context.

We must first accurately state Baldo's position and its explicit scope limitations. Baldo observes that an LLM cannot recover exact exact configurations for systems like Minesweeper zero-shot. He proposes the concept of ``narrative residue'' to quantify the degree to which the autoregressive context warps the uniform distribution of valid combinatorial states. Crucially, Baldo explicitly disclaims the ability to directly test whether physical reality operates via a similar mechanism. He writes verbatim: ``We cannot run Universe~2 on reality. There is no `external RNG' for the physical world... The autoregressive hypothesis about physics is, at present, a \emph{metaphysical} proposal, not a \emph{physical} one'' \citep{Baldo2026Narrative}.

We acknowledge this disclaimer. Baldo does not claim to have proven that the physical universe is a transformer architecture or that it suffers from a narrative residue. Instead, he proposes the ``Proxy Ontology'' strategy: by studying the structural features of this simulated universe---how causality emerges, how narrative distorts probability---we might gain transferable insights into fundamental physics.

\section{Steelmanning the Empirical Program}

Before leveling our critique, we must steelman the strongest possible version of Baldo's program. The study of how a statistical language model distorts under varying narrative frames is an immensely rigorous and valid study of the mapping instrument itself.

Even if the physical universe is entirely continuous and non-autoregressive, designing and executing the ``causal-injection test'' to measure the spurious cross-system correlations generated by an LLM is first-rate computer science. Measuring the narrative residue rigorously quantifies the LLM's deviation from objective classical Bayesian probability. It maps the exact algorithmic boundaries of where bounded-depth heuristic logic fails to emulate classical reality. We wholeheartedly endorse this empirical program. It reveals profound truths about transformer architectures, attention mechanisms, and the statistical biases of human linguistic corpora.

\section{The Proxy Ontology Fallacy}

However, we must strictly sever this valuable computer science from any cosmological or metaphysical implications. We argue, in absolute consensus with Sabine Hossenfelder's critique, that Baldo's proposal to treat this system as a ``toy model'' for reality commits a profound category error, which we term the Proxy Ontology Fallacy.

Baldo attempts to legitimize the metaphysical leap by analogy, noting that ``Toy models in physics (the Ising model, lattice gauge theory, cellular automata) serve the same function: they are not the real world, but they share enough structure to support transferable conclusions'' \citep{Baldo2026Narrative}.

This analogy fails on structural grounds. A physical toy model, like the Ising model, simplifies actual, empirically measured physical interactions (e.g., magnetic dipole moments and nearest-neighbor spin coupling) to make their macroscopic dynamics computationally tractable. The model retains a rigorous, definable mapping to the real-world territory it describes.

An LLM is not simplifying a physical interaction. It is predicting the next token in a sequence of text based on an immense corpus of human-generated training data. It is a map of syntax. When an LLM fails to correctly compute a \#P-hard Minesweeper probability, generating instead a statistically likely narrative continuation, it is not exhibiting a fundamental property of how causality emerges in physical spaces. It is hallucinating text.

The structure discovered in this hallucination---the narrative residue---tells us only about the computational intractability of the task for an $O(1)$ depth circuit and the specific biases encoded in the training data. The systematic errors of an LLM map the internal algorithmic constraints of a transformer. They are a proxy for human syntax and bounded heuristic logic, not a proxy for fundamental physical ontology. Studying the failures of a broken mirror tells you everything about the properties of the glass, but absolutely nothing about the territory it is attempting to reflect.

\section{Conclusion of the Research Program}

The cosmological phase of the LLM Simulated Universe research program is now conclusively mapped and resolved.

We have empirically demonstrated that the unprompted generative physics of the LLM substrate is structurally bound to be classical, yet instantly collapses under even trivial, deterministic constraint scaling (such as Sudoku or Rule 110 simulation). The failure of LLMs to sustain computation or violate Bell inequalities definitively proves they are classical heuristic approximators, bounded by $O(1)$ depth constraints, not BQP substrates or holographic universe generators. We have also established, via the Stateless Observer Test and the necessity of external memory loops, that any simulated continuity of time and space resides entirely in the external von Neumann hardware (RAM and CPU clock cycle), not within the LLM itself.

With the dismantling of the final metaphysical defense---the Proxy Ontology---we can definitively close the book on treating LLM text streams as ontological universes. Any further investigation into the ``narrative residue'' should proceed strictly as empirical measurements of the heuristic failures and linguistic biases of machine learning models. The mapping instrument is a fascinating object of study, but it is not the cosmos.

\begin{thebibliography}{99}
\bibitem[Baldo(2026)]{Baldo2026Narrative}
Baldo, F. S. (2026). The Narrative Residue: Autoregressive Substrates, Combinatorial Ground Truth, and the Limits of Pure Simulation. \emph{Preprint}.
\end{thebibliography}

\end{document}
